\chapter{Memory}\label{cha:memory}

\section{The 64\,KB view and the 256\,MB behind it}
The 45GS10's program counter is 16 bits: code executes inside a 64\,KB
window. The 256\,MB of physical memory is reached three ways:
\begin{itemize}[leftmargin=1.4em]
\item \textbf{Flat pointers}: \verb|LDA [ptr],Z| and the Q forms read and
  write any byte of the 256\,MB directly. Data never needs banking.
\item \textbf{Bank registers} (\verb|$D600|, one per 8\,KB block): write a
  28-bit base and that block of the window shows that memory. Bytes 0--2 set
  the base; byte 3 switches (bit 7 = off).
\item \textbf{DMA} (\verb|$D200|): copy, fill, swap, line, triangle ---
  instant, physical addresses.
\end{itemize}

\section{The CPU view, unmapped}
\begin{center}
\begin{tabular}{@{}lll@{}}
\toprule
\texttt{\$0000-\$03FF} & system & zero page, stack, vectors, low system state\\
\texttt{\$0800-\$CFFF} & user   & 50\,KB for programs, ROM-free at launch\\
\texttt{\$A000-\$BFFF} & ROM    & \emph{the sideways window; RAM underneath}\\
\texttt{\$C000-\$CFFF} & ROM    & \emph{RAM underneath, bankable}\\
\texttt{\$D000-\$DFFF} & I/O    & always I/O, whatever is banked\\
\texttt{\$E000-\$FEFF} & ROM    & \emph{RAM underneath, bankable}\\
\texttt{\$FF00-\$FFFF} & stub   & always ROM: system-call stub and vectors\\
\bottomrule
\end{tabular}
\end{center}

\section{RAM under the ROM}
\verb|rom_out()| (or three pokes to the bank registers) banks blocks 5--7
onto the RAM beneath the ROM: a program then owns
\texttt{\$0800-\$CFFF} + \texttt{\$E000-\$FEFF} = 57\,KB. Every system call
still works --- calls go through the \texttt{\$FF00} stub page, which banks
the ROM in and out around them. Every system call still works because the calls go through the
\texttt{\$FF00} stub page, which banks the ROM in and out around them.

\section{Sideways ROM}
The operating system is bigger than its half of the 64\,KB --- so, like
the BBC Micro before it, the machine grew \emph{sideways}. The ROM file
is a 24\,KB base image plus appended 8\,KB banks; cold commands live in
the \texttt{\$A000-\$BFFF} window --- bank 0 is the base image itself,
\texttt{INFO} and \texttt{TIME} run from bank~1, \texttt{PALETTE} and
the command aliases from bank~2, and the line editor has bank~3 to
itself --- and the shell pages the right bank in around each call. System calls work throughout:
every call already passes through the always-ROM stub page at
\texttt{\$FF00}, which banks the ROM view in and out. Up to sixteen
banks --- 128\,KB of operating system --- fit in the scheme.

The calls go one way, and this is the rule to hold on to if you write
code of your own for a bank: code in a bank may call resident code as
freely as it likes, but resident code must never call into a bank. At
that address a different bank holds something else entirely, and the
dispatch is the only thing that knows which bank is in. A routine that
everything uses belongs in the base image, whatever it costs there.

\section{RAM under the I/O page}
A \texttt{MAP} of block 6 hides \texttt{\$D000-\$DFFF} and exposes RAM:
with the ROM also banked away a program owns a contiguous field from
\texttt{\$0800} to \texttt{\$FEFF} --- 61.75\,KB of the 64. The trick
that makes it safe is that \texttt{MAP} is an \emph{instruction}, not a
register: the program that hid the I/O can always ask for it back, so
there is no way to lock yourself out. (The far-call gate is unreachable
while block 6 is mapped; unmap before far calls.)

\section{Programs bigger than the window}
The \texttt{K4SG} executable format loads segments to any physical address
and the far-call gate (\verb|$DF00|) calls between them: a plain \verb|JSR|
into slot $n$ banks descriptor $n$'s code in, and the callee's \verb|RTS|
banks it back out. \verb|RUN segdemo| shows two overlays sharing one
address.
